% LaTeX file for resume 
% This file uses the resume document class (res.cls)

\documentclass{res} 
\usepackage{hyperref}
\topmargin = -0.4in
\oddsidemargin=-0.3in
\textwidth=6.5in        %
\setlength{\textheight}{9.5in} % increase text height to fit on 1-page 

\begin{document} 

\name{GAUTAM CHETAN KAMATH\\[12pt]}     % the \\[12pt] adds a blank
				        % line after name      

\address{\bf  ADDRESS\\32 Vassar Street\\Room 32-G628\\Cambridge, MA 02139}
\address{\bf CONTACT \\Website: \url{www.gautamkamath.com} \\Cell: (657) 206-7724\\Email: g@csail.mit.edu\\ Skype ID: hoonose\_me}
                                  
\begin{resume}

\section{EDUCATION}
    \vspace{-0.1in}
    \begin{tabbing}
    \hspace{0.3in}\= \kill % set up two tab positions
    {\bf Massachusetts Institute of Technology}  \\        
    \> Ph.D. candidate, September 2018 (expected) \\
    \> Advisor: Constantinos Daskalakis \\
    \> Electrical Engineering and Computer Science \\
    \end{tabbing}
    \vspace{-0.4in}
    \begin{tabbing}
    \hspace{0.3in}\= \kill % set up two tab positions
    {\bf Massachusetts Institute of Technology}  \\        
    \> S.M., September 2014 \\
    \> Thesis: On Learning and Covering Structured Distributions \\
    \> Advisor: Constantinos Daskalakis \\
    \> Electrical Engineering and Computer Science \\
    \end{tabbing}
    \vspace{-0.4in}
    \begin{tabbing}
    \hspace{0.3in}\= \kill % set up two tab positions
    {\bf Cornell University}  \\        
    \> B.S., summa cum laude, May 2012 \\
    \> Computer Science, Electrical and Computer Engineering
    \end{tabbing}

\section{PUBLICATIONS}
      {\bf Which Distribution Distances are Sublinearly Testable?} \\
      Constantinos Daskalakis, Gautam Kamath, John Wright \\
      Proceedings of the 29th Annual ACM-SIAM Symposium on Discrete Algorithms (SODA 2018)

      {\bf Testing Ising Models} \\
      Constantinos Daskalakis, Nishanth Dikkala, Gautam Kamath \\
      Proceedings of the 29th Annual ACM-SIAM Symposium on Discrete Algorithms (SODA 2018)

      {\bf Efficient and Optimally Robust Learning of High-Dimensional Gaussians} \\
      Ilias Diakonikolas, Gautam Kamath, Daniel M. Kane, Jerry Li, Ankur Moitra, Alistair Stewart \\
      Proceedings of the 29th Annual ACM-SIAM Symposium on Discrete Algorithms (SODA 2018)

      {\bf Concentration of Multilinear Functions of the Ising Model with Applications to Network Data} \\
      Constantinos Daskalakis, Nishanth Dikkala, Gautam Kamath \\
      Advances in Neural Information Processing Systems 30 (NIPS 2017)

      {\bf Being Robust (in High Dimensions) Can Be Practical}  \\
      Ilias Diakonikolas, Gautam Kamath, Daniel M. Kane, Jerry Li, Ankur Moitra, Alistair Stewart \\
      Proceedings of the 34th International Conference on Machine Learning (ICML 2017)

      {\bf Priv'IT: Private and Sample Efficient Identity Testing}  \\
      Bryan Cai, Constantinos Daskalakis, Gautam Kamath \\
      Proceedings of the 34th International Conference on Machine Learning (ICML 2017)
     
      {\bf Robust Estimators in High Dimensions without the Computational Intractability} \\
      Ilias Diakonikolas, Gautam Kamath, Daniel M. Kane, Jerry Li, Ankur Moitra, Alistair Stewart \\
      SIAM Journal on Computing, to appear \\
      Preliminary version in Proceedings of the 57th Annual IEEE Symposium on Foundations of Computer Science (FOCS 2016) \\
      {\bf Invited to the SIAM Journal on Computing Special Issue for FOCS 2016} \\
      {\bf Invited to Highlights of Algorithms 2017 (HALG 2017)}

      {\bf A Size-Free CLT for Poisson Multinomials and its Applications} \\
      Constantinos Daskalakis, Anindya De, Gautam Kamath, Christos Tzamos \\
      Proceedings of the 48th ACM Symposium on Theory of Computing (STOC 2016)
      
      {\bf Optimal Testing for Properties of Distributions} \\
      Jayadev Acharya, Constantinos Daskalakis, Gautam Kamath \\
      Advances in Neural Information Processing Systems 28 (NIPS 2015) \\
      {\bf Spotlight Presentation} 

      {\bf On the Structure, Covering, and Learning of Poisson Multinomial Distributions} \\
      Constantinos Daskalakis, Gautam Kamath, Christos Tzamos \\
      Proceedings of the 56th Annual IEEE Symposium on Foundations of Computer Science (FOCS 2015)

      {\bf A Chasm Between Identity and Equivalence Testing with Conditional Queries} \\
      Jayadev Acharya, Cl\'ement Canonne, Gautam Kamath \\ 
      Theory of Computing, to appear \\
      Preliminary version in Proceedings of the 19th International Workshop on Randomization and Computation (RANDOM 2015)

      {\bf Adaptive Estimation in Weighted Group Testing} \\
      Jayadev Acharya, Cl\'ement Canonne, Gautam Kamath \\ 
      Proceedings of the 2015 IEEE International Symposium on Information Theory (ISIT 2015)
 
      {\bf On Learning and Covering Structured Distributions} \\
      Gautam Kamath \\
      S.M. Thesis, 2014 \\
      Available at \url{www.gautamkamath.com/thesis.pdf} 
     
      {\bf Faster and Sample Near-Optimal Algorithms for Proper Learning Mixtures of Gaussians} \\
      Constantinos Daskalakis, Gautam Kamath \\
      Proceedings of the 27th Annual Conference on Learning Theory (COLT 2014)

      {\bf An Analysis of One-Dimensional Schelling Segregation}  \\
      Christina Brandt, Nicole Immorlica, Gautam Kamath, Robert Kleinberg  \\
      Proceedings of the 44th ACM Symposium on Theory of Computing (STOC 2012)

\section{TALKS}
    \vspace{-0.1in}
    \begin{tabbing}
    \hspace{0.3in}\= \kill 
    {\bf Principled Tools for Modern Statistical Data Science}\\
    \> Boston University Computer Science Seminar, February 2018 \\
    \> McGill University Computer Science Seminar, February 2018 \\
    \> University of Waterloo Computer Science Seminar, February 2018 
    \end{tabbing}
    \vspace{-20pt}
    \begin{tabbing}
    \hspace{0.3in}\= \kill 
    {\bf Which Distribution Distances are Sublinearly Testable?} \\
    \> Conference on Information Sciences and Systems ({\bf Invited}), March 2018 \\
    \> Symposium on Discrete Algorithms, January 2018
    \end{tabbing}
    \vspace{-20pt}
    \begin{tabbing}
    \hspace{0.3in}\= \kill 
    {\bf Statistical Hypothesis Testing in the Modern Age} \\
    \> University of Pennsylvania, December 2017 \\
    \> Boston University Algorithms and Theory Seminar, November 2017 \\
    \> University of Massachusetts Amherst Theory Seminar, October 2017 \\
    \> McMaster Seminar in Computing and Software, October 2017 \\
    \> Cornell Theory Seminar, September 2017 
    \end{tabbing}
    \vspace{-20pt}
    \begin{tabbing}
    \hspace{0.3in}\= \kill 
    {\bf Priv'IT: Private and Sample Efficient Identity Testing} \\
    \> International Conference on Machine Learning, August 2017 \\
    \> Private and Secure Machine Learning, August 2017
    \end{tabbing}
    \vspace{-20pt}
    \begin{tabbing}
    \hspace{0.3in}\= \kill 
    {\bf A Size-Free CLT for Poisson Multinomials and its Applications} \\
    \> Symposium on Theory of Computing, June 2016
    \end{tabbing}
    \vspace{-20pt}
    \begin{tabbing}
    \hspace{0.3in}\= \kill 
    {\bf Robust Estimators in High Dimensions without the Computational Intractability} \\
    \> Cornell Theory Lunch, September 2017 \\
    \> China Theory Week, August 2016 ({\bf Invited})
    \end{tabbing}
    \vspace{-20pt}
    \begin{tabbing}
    \hspace{0.3in}\= \kill 
    {\bf Optimal Testing for Properties of Distributions} \\
    \> Northeastern Theory Seminar, March 2017 ({\bf Invited}) \\
    \> University of Pennsylvania Theory Lunch, September 2016 ({\bf Invited}) \\
    \> MIT Signals, Information, and Algorithms Laboratory Group Meeting, March 2016 \\
    \> University of Massachusetts Boston, February 2016 ({\bf Invited}) \\
    \> Berkeley Theory Lunch, September 2015
    \end{tabbing}
    %\vspace{-20pt}
    %\begin{tabbing}
    %\hspace{0.3in}\= \kill 
    %{\bf A Chasm Between Identity and Equivalence Testing with Conditional Queries} \\
    %\> MIT Theory Lunch, February 2015
    %\end{tabbing}
    \vspace{-20pt}
    \begin{tabbing}
    \hspace{0.3in}\= \kill 
    {\bf Faster and Sample Near-Optimal Algorithms for Proper Learning Mixtures of Gaussians} \\
    \> Conference on Learning Theory, June 2014
    %\> MIT Theory Lunch, May 2014
    \end{tabbing}
    \vspace{-20pt}
    \begin{tabbing}
    \hspace{0.3in}\= \kill 
    {\bf An Analysis of One-Dimensional Schelling Segregation} \\
    \> Interdisciplinary Workshop on Information and Decision in Social Networks, November 2012 \\
    \> Symposium on Theory of Computing, May 2012 ({\bf Winner of Best Student Presentation Award})
    \end{tabbing}
    

\section{TEACHING EXPERIENCE}
   \vspace{-0.1in}	
   \begin{tabbing}
   \hspace{0.2in}\= \hspace{4.75in}\= \kill % set up two tab positions
    {\bf Teaching Assistant} \hspace{0.95in}Massachusetts Institute of Technology \hspace{0.25in} Spring 2015, 2017\\
                       \>6.853: Algorithmic Game Theory and Data Science \> (1 semester) \\
                       \>6.856: Randomized Algorithms  \> (1 semester)
   \end{tabbing}
   \vspace{-0.3in}
   \begin{tabbing}
   \hspace{0.2in}\= \hspace{4.75in}\= \kill % set up two tab positions
    {\bf Teaching Assistant} \hspace{0.95in}Cornell University \hspace{1.45in} Spring 2010 - Spring 2012\\
                       \>CS 1114: Intro to Computing with Matlab and Robotics  \> (2 semesters)\\
                       \>CS 2850: Networks  \> (1 semester)\\
                       \>CS 3110: Data Structures and Functional Programming \> (5 semesters)\\
                       \>CS 4820: Introduction to Algorithms  \> (3 semesters)\\
   \end{tabbing}
   \vspace{-20pt}


\section{HONORS AND AWARDS}
    \vspace{-0.1in}
    \begin{tabbing}
    \hspace{5.10in}\=\kill
    MIT Akamai Presidential Graduate Fellowship\>September 2012 - May 2013 \\
    Symposium on Theory of Computing Best Student Presentation Award\>May 2012 \\
    Cornell Computer Science Prize for Academic Excellence\>May 2012 \\
    Eight time Dean's list at Cornell University\>Fall 2008 - Spring 2012 \\
    Recognized by Cornell CS for outstanding work as TA for CS 3110 and CS 4820 \>Spring 2012 \\
    John G. Pertsch Jr. Prize for second highest GPA in ECE \> Spring 2011 \\
    %Second Place in Cornell Facebook Hackathon (\url{http://tinyurl.com/plannit}) \> Spring 2011 \\
    %Featured on Hackaday for ECE 4760 Project (\url{http://tinyurl.com/gckvocal}) \> Spring 2011 \\
    %Attended the Experience Theory Project at University of Washington\>Summer 2010 \\
    Recognized by Cornell CS for outstanding work as TA for CS 1114\>Spring 2010 \\
    %Fourth place in Cornell CS 3110 SteamKart Bot AI Tournament\>Fall 2009 \\
    Canadian Open Mathematics Challenge Gold Medalist in Central Ontario Region\>Spring 2007
    \end{tabbing}

\section{SERVICE}
    \vspace{-0.1in}
    \hspace{5.10in} \\
    Editor of Property Testing Review (March 2016 - Present)\\
    Editor of MIT Theory of Computation Student Blog (November 2013 - Present)\\
    Co-organizer for the TCS+ online seminar series in Theoretical Computer Science (August 2014 - Present)\\
    Co-organizer of FOCS 2017 Workshop on Frontiers in Distribution Testing (October 2017) \\
    Advisor for Danny Lewin MIT Theory Student Retreat (Fall 2014, 2016, 2017) \\
    Co-organizer of FOCS 2016 Workshop on Orthogonal Polynomials and Applications (October 2016) \\
    Organizer of the Second Annual Sublinear Algorithms and Big Data Day (April 2015) \\
    Organizer of Second Annual Danny Lewin MIT Theory Student Retreat (October 2013) \\
    Cofounder and organizer of MIT Theory Lunch (Fall 2012 - Summer 2013) \\
    External reviewer for: STOC, FOCS, SODA, COLT, ITCS, RANDOM, STACS, ICALP, NIPS, \\
    \phantom{asdf} FnT-TCS, Theory Comput. Syst., JMLR


%\section{LEADERSHIP ACTIVITIES}
%    \vspace{-0.1in}
%    \begin{tabbing}
%    \hspace{5.10in}\=\kill
%    President of CSAIL Student Committee \> Spring 2013 - Present\\
%    Organizer of CSAIL Theory Student Lunch \> Fall 2012 - Fall 2013\\
%    President of Cornell Association of Computer Science Undergraduates \> Fall 2011 - Spring 2012\\
%    Junior-at-large of Cornell Association of Computer Science Undergraduates \> Fall 2010 - Spring 2011\\
%    \end{tabbing} 

\section{REFERENCES}
\vspace{-0.1in}
\hspace{5.10in} \\
Constantinos Daskalakis (costis@csail.mit.edu) \\
Associate Professor, Massachusetts Institute of Technology
\end{resume}
\end{document}
